%======================================================================
\section{논의}

%----------------------------------------------------------------------
\subsection{화살표 종류가 어떤 나중 손님을 줄 세울지 정해요. 어느 걷기도 원시 차수를 넘지 못했어요}

\begin{yoji}
어떤 화살표를 고르느냐는 계약 평판 점수가 무엇을 따라가게 할지를 고르는 일이에요. 고른 화살표를 따라 퍼뜨려도 원시 차수보다 나아지지 않았어요(라벨을 안 뒤에 한 비교예요). 대부분의 스펜더에서 엔도스랭크는 무게를 다시 단 허락 수예요.
\end{yoji}

두 화살표 종류는 서로 다른 것을 재고, 기간 밖 결과가 두 창 모두에서 이것을 보여 줘요. 허락으로 만든 점수는 어느 계약이 나중에 새 허락을 얻을지를, 송금으로 만든 점수는 어느 계약이 나중에 새 송금자를 얻을지를 줄 세워요. 그러니 화살표를 고르는 것은, 계약의 평판 점수가 무엇을 따라가야 하는지를 정하는 일이에요.

고른 화살표를 따라 점수를 퍼뜨려도 어느 쪽에서도 보탠 것이 없었어요. 라벨을 안 뒤에 한 비교에서, AWP는 두 창 모두 송금 수보다, 엔도스랭크는 두 창 모두 허락 수보다 낮았어요. 송금을 받은 적 있는 스펜더로 좁힌 집합에서도 W0에서 같았어요. 수를 이기기 어렵다는 것은 선호적 연결(이미 많이 이어진 점에 새 화살표가 더 붙는 현상)과 링크 예측의 관행을 생각하면 놀랍지 않아요~\cite{barabasi1999scaling,libennowell2007link}. 그런데도 온체인 평판 연구는 이것을 거의 시험하지 않아요. 걷기를 다른 걷기와, 또는 자기 화살표로 만든 통계와만 견줘요~\cite{do2023awp}.

닫힌 식은 대부분의 스펜더에서 엔도스랭크가 무엇을 계산하는지 알려 줘요. 허락 수인데, 허락마다 얼마나 최근이고 얼마나 큰지로 크기를 정한 뒤, 그 주인의 허락 무게 합으로 나눈 것이에요. 이 정규화는 주인들이 다른 계약에도 많이 허락하는 스펜더를 낮춰요. 이 정규화 때문인지, 아니면 시간 무게나 금액 무게 때문인지, 무엇이 엔도스랭크와 나중 허락의 일치를 떨어뜨렸는지는 따로 떼어 보지 않았어요.

활동 재시작 변형은 다시 출발 몫을 허락을 가장 많이 주는 주인에게 보내요. 이 변형은 두 창 모두에서 나중의 허락 쌍을 엔도스랭크보다 잘 줄 세웠지만, 여전히 수보다는 아래였어요(표~\ref{tab3}). 그러니 이 비교들에서(모두 라벨을 안 뒤에 한 비교예요) 활발한 주인에서의 다시 출발은 수와의 간격을 좁혔지만 메우지 못했어요. AWP는 논문에 발표된 다시 출발 규칙으로만 돌렸고, 역시 모자랐어요.

걷기가 무게를 다시 단 수보다 더 무언가인 곳에서 퍼뜨림이 수에 무엇을 보태는지 알려면, 두 걸음 이상 바깥까지 추출한 표본과, 아직 추출하지 않은 창에서의 등록한 시험이 필요해요. 그때는 걷기가 무게 단 수로 줄어드는 깊이 1인 스펜더와, 퍼뜨림이 중요할 수 있는 더 깊은 스펜더를 나눠 보고해야 해요.

%----------------------------------------------------------------------
\subsection{트레이더 비교가 보여 주는 것}

\begin{yoji}
재분석에서 허락 수는 주 결과(청산 피하기)에서 송금 수와 갈렸어요. 하지만 이 앞섬은 등록한 시험의 결과가 아니고, 허락을 받은 트레이더 37개와 다른 33개에 기대요. 허락이 신용이나 거래 실력을 잰다는 근거는 없어요.
\end{yoji}

트레이더 결과는 거래소가 기록하고, 어느 그래프도 기록하지 않아요. 이 자료에서 어느 그래프로도 만들지 않은 결과는 이것뿐이에요. 다만 동결 전의 GMX 지급은 송금 그래프에 나타날 수 있어요.

주 결과에서, 이 재분석의 허락 수는 송금 수와 갈렸어요.
\begin{itemize}
\item W1에서 허락 수는 나중의 청산 피하기를 $\tau=0.394$ [0.265, 0.518]로, 송금 수는 0.253 [0.117, 0.379]로 줄 세웠어요. 둘의 차이 P는 97.5\% 구간이 0보다 위였고, 닫은 횟수로 층화한 P도 95\% 구간이 0보다 위였어요.
\item 보조 창 W0은 라벨 기간이 W1의 점수 기간 안에 들어 있어요. 여기서 P는 양수였지만 97.5\% 구간이 0 아래까지 내려왔어요. R1은 W0에서 양수인 추정값만 요구했어요.
\end{itemize}
이 결과에는 두 가지 테두리가 있어요.
\begin{itemize}
\item 점수(달러 금액, 다섯 토큰, 동결 직전 12개월의 점수 기간)는 등록 분석의 결과를 안 뒤에 정의했고, 등록 분석에서는 W1의 P의 97.5\% 구간이 0 아래까지 내려왔어요. 그러니 이 앞섬은 재분석의 발견이지 등록한 시험의 발견이 아니에요.
\item 이 앞섬은 허락을 받은 트레이더 37개와 다른 트레이더 33개에 기대고, 이들은 모두 계약이에요.
\end{itemize}
페이지랭크 쌍은 어느 분석의 어느 창에서도 판가름 나지 않았어요(표~\ref{tab6}). 거래 성공을 나타내는 결과에서는, 허락 쪽 점수와 송금 쪽 점수의 차이가 판가름 난 곳마다 송금 쪽이 앞섰어요. 그리고 어느 점수도 두 창 가운데 어디서도 이익을 낸 비율을 0보다 위의 구간으로 줄 세우지 못했어요.

규칙을 넘어서 보이는 모양은 묘사일 뿐이에요. 허락을 받은 계약은 두 창 모두에서 하나도 청산되지 않았어요. 이 모양은 수십 개의 계약에 기대고, 누가 허락했는지보다 그 계약이 무엇인지를 반영할 수 있어요. 또한 거래 장소가 동결 전에 이미 아는 것, 예를 들어 계약 자신의 청산 기록, 닫은 횟수, 레버리지와 견주지 않았어요. 그러니 과거 행동을 넘어서는 값을 보여 주지 않아요. 그리고 여기의 어떤 결과도 허락이 신용도, 부도 위험, 거래 실력을 잰다고 말해 주지 않아요. 증거금 거래소나 대출 서비스가 허락 수로 계약 상대를 거를 수 있을지는 열린 질문이에요. 그러려면 아직 추출하지 않은 창에서, 더 많은 계약 트레이더로, 과거 행동을 통제한 등록한 시험이 필요해요.

%----------------------------------------------------------------------
\subsection{온체인 평판을 평가하는 방법에 주는 시사점}

\begin{yoji}
걷기 옆에는 늘 원시 차수를 보고하고, 점수는 라벨보다 먼저 동결하고, 어느 화살표로도 만들지 않은 결과로 견주고, 비교마다 언제 정했는지 밝히라는 네 가지 권고예요.
\end{yoji}

이 연구는 온체인 평판 점수를 평가할 때의 네 가지 관행을 권해요.
\begin{enumerate}
\item \textbf{걷기마다 그 옆에 원시 차수를 보고해요.} 여기서 어느 걷기도 기간 밖에서 자기 원시 차수를 이기지 못했어요. 자기 화살표로 만든 통계와만 견준 걷기~\cite{do2023awp}로는 이것을 보일 수 없었을 거예요.
\item \textbf{어떤 라벨을 세기 전에 점수를 동결해요.}
\item \textbf{어느 쪽으로도 만들지 않은 결과에서 화살표 종류를 견줘요.} 각 화살표 종류는 설계상 자기 구성개념의 라벨에서 앞서기 때문이에요~\cite{cronbach1955construct}. 그리고 그 결과가 그래프와 얼마나 떨어져 있는지도 확인해요.
\item \textbf{비교마다 자료를 기준으로 언제 정했는지 밝히고, 결과를 안 뒤에 정의한 분석은 등록한 분석 옆에 둬요}~\cite{nosek2018preregistration}. 여기서 두 분석은 R2와 R3에서 같았어요. 하지만 청산 피하기에서 그냥 센 수의 수준에서 앞선 것은 나중에 정의한 분석(달러 금액, 다섯 토큰, 동결 직전 12개월의 점수 기간)에서만 성립하고, AWP와 활동 변형이 자기 수보다 앞선 것은 등록 분석에서만 나타나요.
\end{enumerate}

%----------------------------------------------------------------------
\subsection{실제 쓰임}

\begin{yoji}
지갑은 주인에게 “이 계약에 양수의 허락을 둔 서로 다른 주인 수”를 보여 줄 수 있어요. 하지만 그 수는 얼마나 널리 허락받았는지를 알려 줄 뿐, 안전하다는 뜻이 아니에요. 그리고 값싸게 부풀릴 수 있어요.
\end{yoji}

어떤 계약에 허락할지 정하는 주인에게, 이미 그 계약에 양수의 허락을 둔 서로 다른 주인 수를 보여 줄 수 있어요. 지갑 프로그램이 계산할 수 있는 공개된 수예요. 기간 밖에서 엔도스랭크는, 균등 재시작이든 활동 재시작이든, 나중의 새 허락 쌍을 이 수만큼 잘 줄 세우지 못했어요.

이런 수는 계약이 이미 얼마나 널리 허락받았는지를 주인에게 알려 줄 뿐, 그 계약이 허락으로 무엇을 할지는 알려 주지 않아요.
\begin{itemize}
\item 피싱 계약도 정직한 라우터처럼 허락을 모아요.
\item 새 계약은 코드가 어떻든 낮은 점수를 받아요.
\item 수가 낮으면 더 살펴볼 까닭이 되지만, 수가 높다고 서명할 까닭이 되지는 않아요.
\end{itemize}
서로 다른 주인을 그냥 센 수는 부풀리기도 싸요. 새 주소 $N$개에서 허락하면, 거래 수수료 $N$번에 토큰은 하나도 묶지 않고 수를 $N$만큼 올릴 수 있어요. 그리고 고르게 다시 출발하는 페이지랭크는 서로 허락해 주는 주소 농장에 열려 있어요~\cite{gyongyi2005alliances,cheng2005sybilproof}. 짝 원고가 그 드러남의 크기를 재요.

어느 계약이 새 송금자를 얻을지 줄 세우는 데서도, 서로 다른 보낸 쪽의 수가 두 창 모두에서 AWP보다 잘 줄 세웠어요(표~\ref{tab3}, 라벨을 안 뒤에 한 비교). 이 점수들은 어느 것도 보안 감사나 신용 점수가 아니고, 담보를 대신하지 않아요~\cite{gudgeon2020crisis}.

%----------------------------------------------------------------------
\subsection{한계}

\begin{yoji}
체인 하나, 12개월 점수 기간, 석 달 라벨 기간 두 개에서 나온 증거예요. 수정 분석(달러 금액, 다섯 토큰, 동결 직전 12개월의 점수 기간)은 등록 분석의 결과를 안 뒤에 정의했고, 트레이더 표본은 작고, 구간은 아마 너무 좁아요.
\end{yoji}

\begin{itemize}
\item \textbf{범위.} 증거는 체인 하나, 12개월 점수 기간, 석 달 라벨 기간 두 개에서 나와요. 다른 롤업과 레이어 1 네트워크는 새 자료가 필요해요~\cite{werner2022sok}.
\item \textbf{토큰.} 그래프는 다섯 토큰만 담아요. 그래서 주인들이 다른 토큰으로만 허락하는 계약은 코호트 밖에 있고, 이더 자체는 WETH로만 들어와요.
\item \textbf{오래된 허락.} 점수 기간보다 오래된 허락은 아직 살아 있어도 읽지 않아요. 그래서 주인들이 오래전에 허락하고 다시 허락하지 않은 계약은 주인이 더 적은 것처럼 점수를 받아요.
\item \textbf{에고 네트워크.} 두 그래프 모두 계약 코호트 하나의 에고 네트워크예요. 그래서 주인의 나가는 무게, 송금자의 활발함, 걷기가 갈 수 있는 경로는 체인 전체가 아니라 그 네트워크의 것이에요. 퍼뜨림에 관한 발견은 보증 경로가 더 긴 그래프로 그대로 옮겨 가지 않아요.
\item \textbf{수정 분석.} 수정 분석은 등록 분석의 결과를 안 뒤에 정의했어요. 그래서 다섯 토큰, 달러 가격의 출처, 기울기, 상한, 동결 직전 12개월의 점수 기간을 모두 그 결과를 본 사람들이 골랐어요. 상수는 W0 점수 기간에서 왔으므로 어떤 라벨도 들어가지 않았고, 등록 분석도 R2와 R3에서 같은 판정에 이르렀어요. 등록 분석은 R1을 만족하지 않았으므로, 그냥 센 수의 수준에서 앞선 것은 나중에 정의한 분석에 기대요.
\item \textbf{AWP의 매개변수.} AWP의 논문이 값을 주지 않아서, 저자들이 골랐어요.
\item \textbf{스펜더 라벨.} 스펜더 라벨은 새로 들어온 것을 센 드문 수라서, 지금의 상대 수를 따르는 점수에 유리해요.
\item \textbf{구간.} 모든 구간은 계약을 서로 독립으로, 그래프를 고정된 것으로 다뤄요. 실제로는 스펜더들이 주인을 나눠 가지므로, 구간은 아마 너무 좁아요. 스펜더 분석은 다시 뽑기를 400번만 썼고, 여러 번 비교한 것에 대한 보정은 판정 쌍에만 했어요.
\item \textbf{계정 종류.} 계정 종류는 두 라벨 기간이 끝난 뒤의 블록에서 읽었어요. 그 전에 스스로 없어진 계약은 외부 소유 계정으로 세어져 빠져요.
\item \textbf{허락 이벤트.} \texttt{Approval} 이벤트는 허락을 줄 때뿐 아니라 쓰고 남은 양을 기록할 때도 나와요. 그리고 Permit2 같은 관리자 안의 허락은 보이지 않아요.
\item \textbf{트레이더.} 트레이더 코호트는 작고, 트레이더 라벨은 일부 닫기도 세요. 청산을 움직이는 시장 상황은 재지 않았어요.
\item \textbf{부도와 인과.} 담보 없는 대출의 부도 라벨은 구할 수 없었고, 인과를 주장하지 않아요.
\end{itemize}

%======================================================================
\section{결론}

\begin{yoji}
허락 화살표와 송금 화살표는 같은 계약에 대해 서로 다른 정보를 담아요. 라벨을 안 뒤에 한 비교에서 어느 걷기도 원시 차수를 넘지 못했어요. 다음 단계는 그냥 센 수의 수준에서 앞선 것을 새 창에서 등록해 시험하는 거예요.
\end{yoji}

허락 화살표와 송금 화살표는 같은 스마트 계약에 대해 서로 다른 정보를 담아요. 아비트럼 원에서 마지막 허락으로 만든 점수는 어느 계약이 새 허락을 얻을지를, 송금으로 만든 점수는 어느 계약이 새 송금자를 얻을지를 줄 세웠어요.

어느 걷기도 자기가 다듬는 원시 차수보다 나아지지 않았어요. 라벨을 안 뒤에 한 비교에서, AWP는 나중에 새 송금자를 얻은 계약들을 두 창 모두 송금 입차수만큼 높은 자리에 놓지 못했어요. 엔도스랭크도 나중에 새 허락 쌍을 얻은 계약들을 두 창 모두 허락한 주인 수만큼 높은 자리에 놓지 못했고, 균등 재시작이든 활동 재시작이든 같았어요. 닫힌 식이 보여 주듯, 대부분의 스펜더에서 엔도스랭크는 무게를 다시 단 허락 수예요.

어느 그래프에도 들어 있지 않은 결과, 곧 계약 트레이더의 나중 청산 피하기에서, 허락 수는 주 창에서 송금 수보다 트레이더를 더 잘 줄 세웠어요. 이것은 그냥 센 수의 수준에서 앞서는지를 따지는 등록 규칙을 이 재분석에서는 만족하지만, 등록 분석에서는 만족하지 않았어요. 페이지랭크 버전은 어느 분석에서도 판가름 나지 않았어요.

다음 단계는 아직 추출하지 않은 창에서, 그 앞섬을 등록한 시험으로 다시 따져 보는 거예요. 더 많은 계약 트레이더를 모으고, 그들의 과거 청산을 통제해야 해요. 걷기를 더 시험한다면, 깊이 1인 스펜더와 더 깊은 스펜더를 나눠 보고해야 해요.
